% !TEX root = main.tex
\documentclass[main.tex]{subfiles}
\begin{document}

\section{Concavity of $\jointdual$}
\label{sec:appendix-concavity}

We now show that our objective $\jointdual$ is concave, using the facet-based structure of its Hessian derived in Appendix~\ref{sec:appendix-joint}.

For conciseness, we stack the two multipliers of a site into $\mult_i = [\shift_i, \weight_i] \in \R^{d+1}$.
We then define the affine field
\begin{equation}
  \psi_i(\xx) = 2\,\shift_i\tr (\xx - \tcentroid_i) + \weight_i = \mult_i\tr \phi_i(\xx) ,
  \qquad
  \phi_i(\xx) = \begin{bmatrix} 2(\xx - \tcentroid_i) \\ 1 \end{bmatrix} ,
  \label{eq:appendix-potential}
\end{equation}
through which $\mult_i$ enters the shifted power distance, since, by \eqref{eq:complete-square}, we have $\smash{\norm{\xx \!-\! \qq_i}^2 \!-\! \psi_i(\xx) \!=\! \norm{\xx \!-\! \dual_i}^2 \!-\! \combweight_i}$.
Given that $\psi_i$ is linear in $\mult_i$, a perturbation $\delta\mult_i \!=\! [\delta\shift_i, \delta\weight_i]$ changes it by
\begin{equation}
  \delta\psi_i(\xx) = \delta\mult_i\tr \phi_i(\xx) = 2\,\delta\shift_i\tr (\xx - \tcentroid_i) + \delta\weight_i.
  \label{eq:appendix-dpsi}
\end{equation}

With this notation, the Hessian blocks \eqref{eq:appendix-blocks-off}--\eqref{eq:appendix-blocks-diag} of $\jointdual$ are the entries of the facet-based Laplacian assembled from $\phi_i$:
\begin{subequations}
  \label{eq:appendix-block}
  \begin{empheq}[left={\empheqlbrace\,}]{align}
    \partial^2_{\mult_i \mult_j} \jointdual
      &= \frac{1}{2}\frac{1}{\norm{\dual_i - \dual_j}} \int_{\face_{ij}} \phi_i(\xx)\, \phi_j(\xx)\tr \, dA,
    \label{eq:appendix-block-off}
    \\[2mm]
    \partial^2_{\mult_i \mult_i} \jointdual
      &= -\!\sum_{j \sim i} \frac{1}{2}\frac{1}{\norm{\dual_i - \dual_j}} \int_{\face_{ij}} \phi_i(\xx)\, \phi_i(\xx)\tr \, dA,
    \label{eq:appendix-block-diag}
  \end{empheq}
\end{subequations}
whose shift--shift part is the block \eqref{eq:centroid-hessian}.

Contracting the Hessian with a perturbation $\smash{\delta\mult = \{\delta\mult_i\}_{i=1}^n}$ and pairing the terms of each facet gives
\begin{equation}
  \delta\mult\tr \nabla^2 \jointdual \, \delta\mult
    =
    -\sum_{\{i,j\}}\!\! \frac{1}{2} \frac{1}{\norm{\dual_i \!-\! \dual_j}} \int_{\face_{ij}}\!\! \bigl( \delta\psi_i(\xx) \!-\! \delta\psi_j(\xx) \bigr)^2 dA
    \le 0 ,
  \label{eq:appendix-nsd}
\end{equation}
where the sum runs over pairs of neighbors.
Hence, $\nabla^2 \jointdual \preceq 0$, and $\jointdual$ is concave.
Note that restricting \eqref{eq:appendix-nsd} to $\delta\weight \!=\! 0$, for which $\delta\psi_i(\xx) \!=\! 2\,\delta\shift_i\tr (\xx - \tcentroid_i)$, shows that, with the weights fixed, $\jointdual$ remains concave with respect to the shifts.

\end{document}
